\section{Game Objects}

Setelah memahami jendela Unity Editor (Bab III), bab ini membahas entitas dasar: GameObject dan Component. \unity{GameObject} adalah entitas dasar dalam Unity \cite{unityManual}. Setiap objek di Hierarchy adalah GameObject. GameObject kosong hanya memiliki \class{Transform}; objek 3D (Cube, Sphere, Plane) memiliki Transform plus Mesh Filter dan Mesh Renderer. Mesh Filter menyimpan data geometri; Mesh Renderer menggambar geometri ke layar.

\subsection{Transform}

Transform mengatur posisi, rotasi, dan skala dalam ruang 3D. Koordinat menggunakan sistem tangan kanan: X kanan, Y atas, Z masuk ke layar.

\begin{table}[htbp]
\centering
\begin{tabular}{|l|p{6cm}|}
\hline
\textbf{Properti} & \textbf{Deskripsi} \\
\hline
Position & Koordinat X, Y, Z dalam ruang world. Origin (0,0,0) di pusat scene \\
\hline
Rotation & Rotasi dalam derajat (Euler). Urutan: X, Y, Z \\
\hline
Scale & Ukuran relatif. (1,1,1) = normal. (2,2,2) = dua kali lipat \\
\hline
\end{tabular}
\caption{Properti Transform}
\end{table}

\subsection{Parent-Child}

GameObject dapat menjadi child dari GameObject lain. Transform child relatif terhadap parent. Memindahkan parent akan memindahkan semua child. Berguna untuk mengelompokkan objek (misal: semua collectible dalam satu parent).

\subsection{Primitive 3D}

Unity menyediakan primitive: Cube, Sphere, Cylinder, Plane, Capsule, Quad. Untuk Roll a Ball, kita gunakan Plane (lantai), Sphere (player), dan Cube (tembok arena).
